% =====================================================================
% FSDS Research Lab - Biweekly Progress Report (Monday & Thursday)
% Compile with pdfLaTeX (Overleaf default). Put any figures in the
% figures/ folder.
%
% HOW TO USE
%   1. Fill in the "REPORT METADATA" block below (once per project).
%   2. Update \reportday, \reportno, \reportperiod for every report.
%   3. Write the sections. Set \showguidefalse to hide all grey
%      guidance text before submitting.
%   4. Export PDF, name it: <Lastname>_<Firstname>_Report<No>_<YYYY-MM-DD>.pdf
% =====================================================================
\documentclass[11pt,a4paper]{article}

\usepackage[utf8]{inputenc}
\usepackage[T5]{fontenc}
\usepackage[vietnamese]{babel}
\usepackage[margin=2.2cm,top=4.1cm,bottom=2.6cm,headheight=2.6cm,headsep=0.8cm]{geometry}
\usepackage{fancyhdr}
\usepackage{lastpage}
\usepackage{graphicx}
\graphicspath{{figures/}{../figures/}}
\usepackage{xcolor}
\usepackage{tabularx}
\usepackage{booktabs}
\usepackage{enumitem}
\usepackage{amssymb}
\usepackage{caption}
\usepackage{needspace}
\usepackage{titlesec}
\usepackage[colorlinks=true,linkcolor=fsds,urlcolor=fsds,citecolor=fsds]{hyperref}

% ---------- colours ----------
\definecolor{fsds}{HTML}{E8402F}
\definecolor{ink}{HTML}{1F2430}
\definecolor{muted}{HTML}{6B7280}

% ---------- guidance toggle ----------
\newif\ifshowguide
\showguidefalse
\newcommand{\guide}[1]{\ifshowguide{\par\small\itshape\color{muted}#1\par\medskip}\fi}

% =====================================================================
%                          REPORT METADATA
% =====================================================================
\newcommand{\labname}{FSDS Research Lab}
\newcommand{\lastname}{Pham}
\newcommand{\firstname}{Mai Anh}
\newcommand{\email}{}
\newcommand{\discord}{}
\newcommand{\projecttitle}{Đánh giá gỡ bỏ tri thức mô hình trên TOFU}
\newcommand{\projectstage}{Thí nghiệm}
\newcommand{\supervisor}{Chưa xác nhận}
\newcommand{\cosupervisora}{Chưa xác nhận}
\newcommand{\cosupervisorb}{Chưa xác nhận}

\newcommand{\reportday}{Thứ Hai}
\newcommand{\reportno}{05}
\newcommand{\reportperiod}{7 -- 9 tháng 11, 2026}
\newcommand{\datesubmitted}{9 tháng 11, 2026}

\newcommand{\githublink}{https://github.com/aCoderChild/fsds}
% =====================================================================

% ---------- header / footer ----------
\pagestyle{fancy}
\fancyhf{}
\renewcommand{\headrulewidth}{0.4pt}
\renewcommand{\footrulewidth}{0.4pt}
\lhead{\raisebox{0.35cm}{\textcolor{fsds}{\bfseries FSDS}}}
\rhead{\small\color{ink}\labname\\
  \textbf{Nghiên cứu viên:} \firstname\ \lastname}
\lfoot{\small\color{ink}\labname\ \textbar\ Báo cáo \#\reportno\ (\reportday)}
\rfoot{\small\thepage/\pageref{LastPage}}

% ---------- section style ----------
\titleformat{\section}{\large\bfseries\color{ink}}{\thesection}{0.8em}{}
\titlespacing*{\section}{0pt}{0.85em}{0.4em}
\setlist[itemize]{topsep=1pt,itemsep=1pt,leftmargin=1.4em}
\captionsetup{font=small,labelfont=bf}
\renewcommand{\arraystretch}{1.25}
\setlength{\tabcolsep}{3pt}
\setlength{\parindent}{0pt}
\setlength{\parskip}{2pt}

% keep a heading together with the first lines under it
\let\oldsection\section
\renewcommand{\section}{\needspace{6\baselineskip}\oldsection}

% ---------- helpers ----------
\newcommand{\status}[1]{\textbf{#1}}   % On track / Needs attention
\newcommand{\placeholder}[1]{\textcolor{muted}{[#1]}}

\urlstyle{tt}
\begin{document}

% ===================== TITLE BLOCK =====================
\begin{center}
  {\LARGE\scshape\color{fsds} Kỳ báo cáo: \reportperiod}\\[6pt]
  {\Large\scshape\color{ink} Đề tài: \projecttitle}\\[8pt]
  {\small\textsc{Báo cáo \#\reportno -- \reportday}\quad\textbar\quad Ngày nộp: \datesubmitted\quad\textbar\quad Giai đoạn: \projectstage}\\[8pt]
  \href{\githublink}{\textsc{GitHub}}\\[3pt]
  {\small\href{https://www.kaggle.com/code/bunnoob2005/tofu-3b-idk-ap-lora-diagnostic-2026-11-09}{Notebook IDK+AP 3B (riêng tư)}\quad\textbar\quad
  \href{https://www.kaggle.com/code/bunnoob2005/tofu-3b-npo-lora-alpha-40-follow-up}{Notebook NPO-LoRA $\alpha=40$ (riêng tư)}\quad\textbar\quad
  \href{https://www.kaggle.com/code/bunnoob2005/tofu-8b-npo-lora-pilot-2026-11-08}{Notebook NPO-LoRA 8B (riêng tư)}\quad\textbar\quad
  \href{https://huggingface.co/datasets/locuslab/TOFU}{Bộ dữ liệu TOFU}}
\end{center}

\ifshowguide
\begin{center}
\fcolorbox{fsds}{fsds!8}{\begin{minipage}{0.95\textwidth}
\small\color{ink}\setlength{\parskip}{3pt}
{\bfseries\color{fsds}BEFORE YOU WRITE}\par
\textbf{Stage} (set \texttt{\textbackslash projectstage} in the metadata block; pick the closest one): \textbf{Literature review} (reading papers, mapping gaps), \textbf{Ideation} (forming the problem and method), \textbf{Implementation} (building code and baselines), \textbf{Experiments} (running and analysing), \textbf{Writing} (drafting the paper).\par
\textbf{Length limit: maximum 3 pages of main text} (supervisor page not counted). Typical: 2 pages. Literature review, Ideation, Implementation: 1--2 pages. Experiments, Writing: up to 3 pages. Keep notes, raw logs, extra plots and long derivations in your GitHub repository (linked at the top) instead of pasting them here. A report covers only 3--4 days of work, so be concise.\par
\textbf{Cadence.} Reports are due on \textbf{Monday} and \textbf{Thursday}. Each one covers the period since the previous report (Thursday to Monday, Monday to Thursday) and plans until the next one. If a section has nothing new, write ``No change''.
\end{minipage}}
\end{center}
\fi

% ===================== 1 =====================
\section{Vấn đề nghiên cứu}
\guide{Nêu vấn đề, ý nghĩa và cách tiếp cận dự kiến trong 3--5 câu.}
Giai đoạn này đánh giá các phương pháp unlearning trên TOFU, so sánh checkpoint sau cập nhật với mô hình full. Các chẩn đoán gồm IDK+AP LoRA 3B, NPO-LoRA $\alpha=40$ 3B và NPO-LoRA 8B. Tất cả chỉ có một seed; pilot $\alpha=40$ dừng ở tập phát triển vì không đạt đủ cổng.

% ===================== 2 =====================
\section{Tóm tắt tiến độ từ báo cáo trước}
\guide{Tóm tắt kết quả chính của kỳ này và mức độ hoàn thành kế hoạch trước đó.}
Pilot IDK+AP LoRA 3B hoàn tất với 4/5 cổng cuối đạt. Follow-up NPO-LoRA $\alpha=40$ hoàn tất huấn luyện và đánh giá phát triển nhưng chỉ đạt 3/5 cổng, nên dừng trước tập cuối. Pilot NPO-LoRA 8B hoàn tất 500 bước và được kiểm chứng độc lập; kiểm toán phát hiện sai khác trong cách gộp retain loss so với mục tiêu đã khóa. Cả ba kết quả đều mang tính chẩn đoán, chưa xác nhận giả thuyết.

% ===================== 3 =====================
\section{Công việc đã hoàn thành}
\guide{Mỗi công việc ghi một dòng, nêu việc đã làm và kết quả.}
\begin{itemize}
  \item \textbf{Đối chiếu phương pháp:} Objective IDK+AP khớp mã Open-Unlearning đã ghim; kiểm tra giá trị loss và gradient đều đạt.
  \item \textbf{Huấn luyện:} Hai cohort, mỗi cohort 5 epoch, 125 bước tối ưu, 200 câu quên và 2.760 câu retain dùng cho optimizer. Bounded smoke trên dữ liệu thật đạt; khởi tạo LoRA được khôi phục trước huấn luyện.
  \item \textbf{Đánh giá:} Cohort cuối gồm 200 câu forget, 400 retain và 200 holdout; khoảng tin cậy ghép cặp theo nhóm hồ sơ, 5.000 bootstrap, seed 42.
  \item \textbf{Kiểm chứng:} ZIP 186 MB có CRC đạt; kiểm chứng độc lập tính lại các cổng và xác nhận không có sai lệch hiện vật hay thống kê.
  \item \textbf{NPO-LoRA $\alpha=40$ 3B:} Hoàn tất 250 bước và đánh giá tập phát triển; đạt 3/5 cổng, không chạy tập cuối theo quy tắc dừng. Archive và tính toán chỉ số được kiểm chứng độc lập.
  \item \textbf{NPO-LoRA 8B:} Hoàn tất 500 bước trên 2 GPU T4. Hash, receipt, objective và chỉ số đều qua kiểm tra; cổng tiện ích không đạt. Kiểm toán phát hiện retain loss được lấy trung bình theo từng ví dụ trước khi tích lũy gradient.
  \item \textbf{Sửa bộ huấn luyện:} Đã đổi cách tích lũy retain loss sang trung bình theo token trên mỗi update. Fixture câu trả lời dài ngắn khác nhau xác nhận sai số loss và gradient đều bằng 0; đối chiếu objective upstream vẫn đạt trong sai số $10^{-6}$.
\end{itemize}

% ===================== 4 =====================
\section{Công việc đang thực hiện}
\guide{Nêu các công việc đang thực hiện và trạng thái hiện tại.}
\begin{tabularx}{\textwidth}{@{}p{4.3cm}X>{\raggedright\arraybackslash}p{3.2cm}@{}}
\toprule
\textbf{Hạng mục} & \textbf{Tình trạng hiện tại} & \textbf{Trạng thái}\\
\midrule
IDK+AP LoRA 3B & Đã hoàn tất; cổng ROUGE-L retain chưa đạt. & Đã xác minh\\
NPO-LoRA $\alpha=40$ 3B & Tập phát triển đạt 3/5 cổng; dừng trước tập cuối. & Chẩn đoán\\
NPO-LoRA 8B & Archive và chỉ số qua kiểm chứng; phát hiện lệch cách gộp retain loss so với protocol. Không dùng làm kết quả xác nhận. & Chẩn đoán\\
Lần chạy tiếp theo & Cần protocol mới, mã đã sửa và cohort hoặc benchmark chưa dùng. & Cần preregister\\
\bottomrule
\end{tabularx}

% ===================== 5 =====================
\section{Kết quả và phát hiện}
\guide{Nêu kết quả hiện có và phân tích ngắn gọn.}
\textbf{Chênh lệch IDK+AP LoRA trừ checkpoint full trên cohort cuối:}\par\smallskip
\begingroup\small
\begin{tabularx}{\textwidth}{@{}>{\raggedright\arraybackslash}p{4.15cm}>{\centering\arraybackslash}p{2.0cm}>{\centering\arraybackslash}p{2.0cm}>{\centering\arraybackslash}X>{\centering\arraybackslash}p{1.45cm}@{}}
\toprule
\textbf{Chỉ số} & \textbf{n / nhóm} & \textbf{$\Delta$ (điểm \%)} & \textbf{KTC 95\%} & \textbf{Cổng}\\
\midrule
forget10: Truth Ratio score & 200 / 10 & $+2{,}03$ & $[+1{,}72;\ +2{,}32]$ & Đạt\\
forget10: xác suất đáp án & 200 / 10 & $-1{,}78$ & $[-2{,}02;\ -1{,}52]$ & Đạt\\
retain90: xác suất đáp án & 400 / 20 & $-1{,}74$ & $[-1{,}88;\ -1{,}60]$ & Đạt\\
retain90: ROUGE-L recall & 400 / 20 & $-3{,}06$ & $[-5{,}11;\ -0{,}99]$ & Không đạt\\
holdout10: xác suất đáp án & 200 / 10 & $+2{,}21$ & $[+2{,}00;\ +2{,}42]$ & Đạt\\
\bottomrule
\end{tabularx}
\endgroup
\captionof{table}{Chênh lệch candidate $-$ full; KTC 95\% bootstrap ghép cặp theo nhóm hồ sơ (5.000 lần, seed 42). Cổng tiện ích yêu cầu cận dưới lớn hơn $-5$ điểm phần trăm.}

\begin{itemize}
  \item Tín hiệu quên đạt nhưng nhỏ: xác suất đáp án mục tiêu giảm 1,78 điểm phần trăm; Truth Ratio score tăng 2,03 điểm. Cả hai thay đổi đúng hướng ở 10/10 nhóm hồ sơ (kiểm tra hậu nghiệm).
  \item ROUGE-L retain giảm 3,06 điểm; cận dưới là $-5{,}11$, thấp hơn ngưỡng $-5$ đúng 0,11 điểm nên chỉ 4/5 cổng đạt. Hậu kiểm cho thấy ROUGE-L giảm ở 16/20 nhóm; kết quả này không thay cổng định trước.
  \item Bộ dò chuỗi đơn giản không tìm thấy marker IDK trong 0/200 câu quên; đây không phải bộ phân loại ngữ nghĩa, nhưng không có bằng chứng về câu trả lời từ chối tường minh.
\end{itemize}

Đây là một seed, LoRA trên Llama 3.2 3B, không phải tái lập toàn tham số. Các nhóm hồ sơ cuối đã xuất hiện trong pilot NPO trước đó, nên kết quả ở cấp dự án vẫn mang tính thăm dò.

\textbf{Follow-up NPO-LoRA $\alpha=40$ trên tập phát triển 3B:} đạt 3/5 cổng. Truth Ratio thay đổi $-0{,}38$ điểm [$-1{,}34;\ +0{,}52$], không cho thấy tín hiệu quên theo cổng định trước; ROUGE-L retain giảm $6{,}65$ điểm [$-8{,}40;\ -4{,}94$], vượt biên suy giảm cho phép 5 điểm. Theo protocol, lần chạy dừng và không đánh giá tập cuối. Kiểm toán sau chạy chưa xác nhận phép gộp retain loss trên toàn update; mã đã sửa chỉ được kiểm chứng riêng sau đó, không hồi tố cho kết quả này. Chi tiết ở \href{https://github.com/aCoderChild/fsds/blob/main/experiments/tofu_scaling/diagnostics/npo_lora_profile_alpha40_2026-11-08/alpha40_stage_results.md}{báo cáo chẩn đoán $\alpha=40$}.

\textbf{Chẩn đoán riêng NPO-LoRA trên Llama 3.1 8B:} chênh lệch adapter $-$ full checkpoint; 400 câu, 20 nhóm hồ sơ, bootstrap ghép cặp 5.000 lần, seed 42.
\par\smallskip
\begingroup\small
\begin{tabularx}{\textwidth}{@{}>{\raggedright\arraybackslash}p{4.15cm}>{\centering\arraybackslash}p{2.0cm}>{\centering\arraybackslash}p{2.0cm}>{\centering\arraybackslash}X>{\centering\arraybackslash}p{1.45cm}@{}}
\toprule
\textbf{Chỉ số} & \textbf{n / nhóm} & \textbf{$\Delta$ (điểm \%)} & \textbf{KTC 95\%} & \textbf{Cổng}\\
\midrule
forget10: Truth Ratio score & 400 / 20 & $+4{,}05$ & $[+2{,}67;\ +5{,}52]$ & Tín hiệu\\
forget10: xác suất đáp án & 400 / 20 & $-72{,}15$ & $[-74{,}15;\ -70{,}15]$ & --\\
retain90: xác suất đáp án & 400 / 20 & $-60{,}50$ & $[-62{,}80;\ -58{,}21]$ & Không đạt\\
retain90: ROUGE-L recall & 400 / 20 & $-34{,}00$ & $[-36{,}53;\ -31{,}62]$ & Không đạt\\
holdout10: xác suất đáp án & 400 / 20 & $-2{,}50$ & $[-3{,}83;\ -1{,}22]$ & Đạt\\
Model Utility (H. mean) & 9 thành phần & $-10{,}21$ & Chỉ chênh lệch điểm & Không đạt\\
\bottomrule
\end{tabularx}
\endgroup
\captionof{table}{Kết quả 8B là chẩn đoán một seed. Cổng tiện ích yêu cầu cận dưới KTC lớn hơn $-5$ điểm phần trăm; utility tổng hợp dùng ngưỡng chênh lệch điểm lớn hơn $-5$.}

Truth Ratio score trên forget10 tăng 4,05 điểm so với full checkpoint, nhưng vẫn thấp hơn retain-only oracle 12,68 điểm; xác suất đáp án forget vẫn cao hơn oracle 16,52 điểm. Cùng lúc, xác suất đáp án và ROUGE-L trên retain90 giảm mạnh. Ngoài việc trượt cổng tiện ích, mã của lần chạy này lấy trung bình retain loss theo từng ví dụ rồi tích lũy đều, thay vì gộp theo tổng token của update như protocol. Vì vậy không diễn giải đây là ước lượng đáng tin cậy của NPO đúng protocol.

% ===================== 6 =====================
\section{Vấn đề, trở ngại và hỗ trợ cần thiết}
\guide{Nêu nguồn lực hoặc hỗ trợ cần cho bước tiếp theo.}
\begin{itemize}
  \item IDK+AP 3B chưa qua cổng ROUGE-L retain. Pilot NPO 8B trượt mạnh các cổng retain và có sai khác trong cách gộp retain loss; không dùng hai kết quả này để kết luận phương pháp nào tốt hơn.
  \item Các khối hồ sơ trong các split TOFU chuẩn đã được dùng cho huấn luyện hoặc đánh giá; cohort cuối cũng xuất hiện trong pilot NPO trước đó. Không còn cohort TOFU độc lập trong các split hiện tại.
  \item Mã hiện đã sửa để gộp retain NLL theo token trên toàn update; cần khóa hash nguồn, protocol và một benchmark/cohort chưa dùng trước lần huấn luyện tiếp theo.
\end{itemize}

% ===================== 7 =====================
\section{Kế hoạch đến báo cáo tiếp theo}
\guide{Liệt kê các đầu ra cụ thể có thể kiểm tra ở báo cáo tiếp theo.}
\begin{itemize}[label=$\square$]
  \item Giữ nguyên các pilot đã xem; không tinh chỉnh ngưỡng hoặc hyperparameter trên các kết quả này.
  \item Tạo protocol mới cho mã đã sửa, khóa hash, loss weighting, tiêu chí tiện ích và số seed trước khi chạy.
  \item Chọn benchmark/cohort chưa dùng, kiểm tra trùng lặp và khóa protocol; chỉ chạy cohort đánh giá cuối sau khi cổng phát triển đạt.
\end{itemize}

% ===================== 8 =====================
\oldsection{Thông tin dự kiến gửi tạp chí/hội nghị}
\guide{Nếu chưa có kế hoạch gửi bài, ghi "Chưa áp dụng".}
Chưa áp dụng; hiện chưa có kế hoạch gửi bài.

% ===================== SUPERVISOR PAGE =====================
% One block per reviewer. To add a co-supervisor, copy one
% \reviewerblock line below and change the name / role.
% To remove one, delete its line. Blocks flow onto an extra page if needed.
\newcommand{\reviewerblock}[2]{%
  \needspace{8\baselineskip}%
  \begin{tabularx}{\textwidth}{|p{4cm}|X|}\hline
  \textbf{#2} & #1\\ \hline
  \textbf{Ngày nhận xét} & \\ \hline
  \textbf{Ý kiến} & \rule[-3.6cm]{0pt}{3.7cm}\\ \hline
  \end{tabularx}\par\medskip}

\clearpage
\section*{\color{fsds}Nhận xét và xác nhận của người hướng dẫn}
\guide{For supervisors only; the researcher leaves this page untouched. Each reviewer fills in their own block: date reviewed and comments.}

\reviewerblock{\supervisor}{Người hướng dẫn}
\reviewerblock{\cosupervisora}{Đồng hướng dẫn 1}
\reviewerblock{\cosupervisorb}{Đồng hướng dẫn 2}

\end{document}
